
\section{Are the Pastors the "Society’s" Servants?}

"Folkebladet" for March 7, 1894.

By "the society" is meant in this connection the organized church body which at its annual meetings decides its common affairs by majority vote.

I see that people have now furthermore begun to call the pastors "the society’s servants," just as in earlier times these were called "the Lord’s servants," "the Word’s servants," "the congregation’s servants." People speak of their "standing in the service of the United Church," so long as they are proclaimers of the Gospel in congregations which belong to this society.

This manner of speech, of which there are more and more examples at least within our society, may naturally be innocent enough, if it rests upon thoughtlessness or linguistic incompetence; but it may also have arisen from a conception of the pastor’s position both crooked and dangerous.

That the pastor is the Lord’s servant expresses his independence of men and his responsibility before God; that the pastor is the congregation’s servant expresses his self-sacrificing labor and suffering for those for whom Christ labored and suffered. Both names are so well grounded in the Word of God and so well understood among all spiritual people that it is surely superfluous to explain them to Christians; and those who would misunderstand these titles of honor would also misunderstand our explanation of them.

It is otherwise, however, with the modern expression: "the society’s servants" and "pastors who stand in the society’s service." These expressions are wholly foreign to Scripture and therefore also foreign to the way of thinking of spiritual people.

It is therefore not superfluous to ask what is meant by these expressions, and with what intention they are used so diligently.

If by "the society’s servants" one meant such persons as seek the society’s good and with the sacrifice of time and strength labor for it, men who set the society’s best above the favor of the society’s majority, then it might do well enough to call the pastors the society’s servants. In that sense Paul says that he has become the servant of all; at the same time he is a mighty man in asserting his independence of all, and rebuking it as a sin if anyone seeks to please men by surrendering anything of the truth.

But if something else is meant by it, namely that the pastors are the organs, instruments, and agents of a higher authority in the congregations, then this designation begins to become dangerous and dubious.

Then the question at once arises: Is the society an authority over the congregations? Is the society a higher power over the congregations? And is it really this higher power whose servants the pastors are, so that, if things went as they ought, the society should rather send its servants to the congregations, and not the congregations choose their pastors?

But by that we have entered into a way of thinking far from being according to Christ, and very well fitted to lower the pastors and rob them of their rightful loftiness and dignity.

For a servant does not know what his master does, and has nothing to do with asking after it. A servant in this sense is one who, when his master says to him, Go! then goes; Come! then comes. He does not ask why, and does not reflect upon it; obedience is his duty, and he shows it without more ado. If such a servant thinks that what is demanded of him is unreasonable, then he can leave the service; therein consists his freedom.

If now the pastors in this sense are the society’s servants, that is, if it is their duty without more ado to carry out the determinations of the majority and persuade their congregations to go in for them, then it is accordingly also the pastors’ duty to leave the society if they are not in every respect agreed with the majority’s plans. If some pastor finds that the society has acted badly and unjustly, then it is not right for him to remind, admonish, rebuke, and convince the society; he is, after all, a servant. Neither is it right for him to remain standing in the society’s service; there is for him only one single way out, namely, to leave the society.

But if one really wishes to proceed according to this view, which now seems to be asserting itself both in the home mission matter and in other ways, then one of two things must happen: either the society will become smaller and smaller, as every opposition is continually driven out; or else the pastors will become men without conviction and conscience. And the more compliant a majority finds the pastors, the more matters it draws under its decision, until at last there is no ecclesiastical question left which is not decided by majority vote.

The pastors ought to be the Lord’s servants and the truth’s undaunted witnesses. Before any others their voice ought to be lifted up like a trumpet to proclaim to the individual and to the society its sins.

And the society, so far from seeking to dampen and stifle this voice either by money or by other means, ought on the contrary to reward with the greatest thankfulness that pastor who corrected and set right an error.

The society is best served by the most self-dependent and independent servants of the Lord, who dare to say of evil that it is evil, even when they must stand as a fortified wall and an iron pillar against Judah’s kings and priests and princes and against the people of the land. But such men cannot bend their conviction concerning right and wrong according to the decree of an accidental majority.

Georg Sverdrup
